\documentclass[french]{article}

\usepackage{babel}
\usepackage{comment}
 
\pagestyle{headings}
\markright{28/6/1995 \hfill Conservatoire National des Arts et
M\'etiers -- Paris \hfill\ } 

\baselineskip 6pt
\parskip 6pt
\onecolumn
\textwidth 16cm
\textheight 24cm
\hoffset=-2cm
\voffset=-2.5cm
\setlength{\parindent}{0.0cm}
\unitlength 1cm
\newtheorem{defin}{D\'efinition}[section]
\newtheorem{exampl}{Exemple}[section]
\def\picwidth{16cm}

\newenvironment{solution}{\nopagebreak\par\vspace*{2mm}\underline{Solution~:}\em\par}{\par\vspace*{2mm}}

%\excludecomment{solution}

\begin{document}
\begin{center}
\Large\bf Processus d'Informatisation et Bases de Donn\'ees A \\ 
    Examen Bases de Donn\'ees\\
\large\bf Mercredi 23 juin 1999
\end{center}

%%%%%%%%%%%%%%%%%%%%%%%%%%%%%%%%%%%%%%%%%%%%%%%%%%%%%%%%%%%%

\vspace*{2cm}


Des \'editeurs se r\'eunissent pour cr\'eer une Base
de Donn\'ees sur leurs publications scientifiques.
Dans de telles publications, plusieurs auteurs se
regroupent pour \'ecrire un livre en se r\'epartissant les chapitres
\`a r\'ediger.
Apr\`es discussion, voici le sch\'ema obtenu:
 
{\tt
\begin{enumerate}
   \item LIVRE (\underline{titreLivre}, ann\'ee,
                          \'editeur,  chiffreAffaire)
   \item CHAPITRE (\underline{titreLivre, titreChapitre}, nbPages)
   \item AUTEUR (\underline{nom}, pr\'enom, ann\'eeNaissance)
   \item REDACTION (\underline{nomAuteur, titreLivre, titreChapitre})
\end{enumerate}
}

Les cl\'es primaires sont soulign\'ees, mais pas les cl\'es \'etrang\`eres.

\section{Cr\'eation du sch\'ema}

\begin{enumerate}
  \item ({\bf 4 pts})
        Donnez les ordres {\tt CREATE TABLE} pour le sch\'ema,
       en sp\'ecifiant soigneusement cl\'es primaires et \'etrang\`eres
      avec la syntaxe SQL2. Le type des donn\'ees est secondaire:
     choisissez ce qui vous semble logique.

\begin{solution}

\begin{verbatim}
CREATE TABLE Livre (titreLivre VARCHAR2 (30),
	            ann�eParution    NUMBER (4),
                    �diteur    VARCHAR2 (30),
	            chiffreAffaire   NUMBER (12,2),
	             PRIMARY KEY (titreLivre)
                    );

CREATE TABLE Chapitre (titreLivre   VARCHAR2 (30),
                       titreChapitre  VARCHAR2 (30),,
    	               nbPages   NUMBER (4),
                       PRIMARY KEY (titreLivre, titreChapitre),
                       FOREIGN KEY (titreLivre) REFERENCES Livre
                         ON DELETE CASCADE
                       );

CREATE TABLE Auteur (nom VARCHAR2 (30),
	             pr�nom   VARCHAR2 (30),
                     ann�eNaissance  NUMBER (4),
	             PRIMARY KEY (nom)
                    );

CREATE TABLE R�daction (nomAuteur VARCHAR2 (30),
                     titreLivre   VARCHAR2 (30),
                     titreChapitre  VARCHAR2 (30),,
               PRIMARY KEY (nomAuteur,titreLivre, titreChapitre),
               FOREIGN KEY (nomAuteur) REFERENCES Auteur
                               ON DELETE CASCADE,
               FOREIGN KEY (titreLivre,titreChapitre) REFERENCES Livre
                               ON DELETE CASCADE
                       );
\end{verbatim}

La table {\tt REDACTION} est un exemple standard d'association n-m.
\end{solution}

   \item ({\bf 1 pt}) Sur quelles cl\'es \'etrang\`eres devrait-on mettre
          l'option {\tt ON DELETE CASCADE}\,?


   \item  ({\bf 2 pts}) Donnez le sch\'ema Entit\'e/Association,
     en Merise ou en UML, correspondant au sch\'ema relationnel 
     de l'\'enonc\'e.
\begin{solution}
Association de composition entre {\tt Livre} et {\tt Chapitre}.
Association {\tt n-n} entre chapitre et auteur.

\end{solution}
   \item ({\bf 2 pts}) Voici les d\'ependances fonctionnelles pour la table
         {\tt LIVRE}\,: \{$titreLivre \to annee; titreLivre \to editeur;
                            editeur \to chiffreAffaire$\}
        Montrer que la table n'est pas en troisi\`eme forme normale et
       proposez une bonne d\'ecomposition.

\begin{solution}
La DF  $editeur \to chiffreAffaire$ pose probl�me puisque 
$editeur$ n'est pas une cl\'e, et $chiffreAffaire$ ne fait pas partie
de la cl\'e. Donc il faut d\'ecomposer en 
{\tt LIVRE (titreLivre, editeur, ann�e)} et 
{\tt EDITEUR (editeur, chiffreAffaire)}. {\tt �diteur} devient
cl\'e \'etrang�re dans la table {\tt Livre}.

\end{solution}
\end{enumerate}

\end{document}

\section{Requ\^etes}

Sur le sch\'ema {\bf propos\'e dans l'\'enonc\'e}, exprimez
en alg\`ebre relationnelle {\bf ou} en SQL les requ\^etes suivantes:

\begin{enumerate}
  \item ({\bf 1 pt}) Titre des livres \'edit\'es chez l'\'editeur 'Vulcain'.
\begin{solution}
SELECT titre FROM livre WHERE editeur='Vulcain'
\end{solution}

  \item  ({\bf 1 pt}) Quel(s) \'editeur(s) a (ou ont) \'edit\'e 
       un livre dont un des chapitres
        a pour titre 'SQL'\,?
\begin{solution}
SELECT editeur 
FROM Livre l, Chapitre c
WHERE l.titreLivre = c.titreLivre
AND  titreChapitre = 'SQL'
\end{solution}

   \item  ({\bf 1 pts}) Nom des auteurs qui ont contribu\'e \`a un chapitre 
      de plus de 50 pages.
\begin{solution}
SELECT nomAuteur
FROM Redaction r, Chapitre c
WHERE  r.titreLivre = c.titreLivre
AND  r.titreChapitre = c.titreChapitre
AND  nbPages > 50
\end{solution}

    \item  ({\bf 2 pts})
      Nom des auteurs 
         qui ont contribu\'e \`a un livre paru chez Vulcain
\begin{solution}
SELECT nomAuteur
FROM  Redaction r, Livre v
WHERE  r.titreLivre = l.titreLivre
AND   editeur = 'Vulcain'
\end{solution}

   \item  ({\bf 2 pts})
       Titre des livres chez Vulcain qui ne contiennent 
              pas de chapitre sur SQL (titreChapitre='SQL')
\begin{solution}
SELECT titreLivre 
FROM    Livre
WHERE   editeur = 'Vulcain'
AND     titreLivre NOT IN (SELECT titreLivre FROM chapitre
                           WHERE titreChapitre='SQL')
\end{solution}

   \item  ({\bf 2 pts}, SQL seulement) Nom des auteurs
      qui ont contribu\'e \`a un livre paru lors qu'ils
      avaient moins de 25 ans.
\begin{solution}
SELECT nom
FROM  Auteur a,Redaction r, Chapitre c, Livre v
WHERE  r.titreLivre = c.titreLivre
AND  r.titreChapitre = c.titreChapitre
AND  l.titreLivre = c.titreLivre
AND   a.nom = r.nomAuteur
AND   anneeParution - anneeNaissance <= 25
\end{solution}

   \item  ({\bf 1 pt}, SQL seulement) Donner, pour chaque \'editeur, le nombre
          de livres qu'il a publi\'e.

\begin{solution}
SELECT editeur, count(*)
FROM Livre
GROUP BY editeur
\end{solution}

   \item  ({\bf 2 pts}, SQL seulement) Donner les titres des livres ayant plus
        de 7 chapitres.
\begin{solution}
SELECT titreLivre
FROM Chapitre
GROUP BY titreLivre
HAVING count(*) >= 7
\end{solution}

   \item  ({\bf 1 pt}, SQL seulement) Donner l'ordre SQL ins\'erant
       l'auteur Ben Duchemin, n\'e en 1963, puis
            l'ordre SQL
          modifiant l'ann\'ee de naissance de Duchemin qui
          devient 1965.

\begin{solution}
INSERT INTO Auteur VALUES ('Duchemin', 'Ben', 1963)
UPDATE Auteur SET anneeNaissance = 1965 WHERE nom = 'Duchemin'
\end{solution}
\end{enumerate}

\begin{center}

\framebox(10,2){{\bf Solution sur le WEB, et au secr\'etariat}}
\end{center}
\end{document}